% quickstart_guide.tex
% A succinct quick-start guide for LLM-experienced developers new to Claude Code
% Build: make quickstart

\documentclass[10pt,oneside]{scrartcl}

% --- Geometry: narrower for a reference guide feel ---
\usepackage{geometry}
\geometry{
    paperwidth=8.5in,
    paperheight=11in,
    left=0.9in,
    right=0.9in,
    top=0.8in,
    bottom=0.8in,
    footskip=0.4in,
    headheight=20pt
}

% --- Engine ---
\usepackage{iftex}
\RequireLuaTeX

% --- Fonts ---
\usepackage{fontspec}

\directlua{
  luaotfload.add_fallback("symbolfallback", {
    "DejaVu Sans:mode=harf;",
  })
}

\setmainfont{Roboto}[
  UprightFont = *-Regular,
  BoldFont = *-Bold,
  ItalicFont = *-Italic,
  BoldItalicFont = *-BoldItalic,
  Scale = 1.0,
  RawFeature={fallback=symbolfallback}
]

\setmonofont{Source Code Pro}[
  Scale = 0.82,
  UprightFont = *-Regular,
  BoldFont = *-Bold,
  ItalicFont = *-It,
  BoldItalicFont = *-BoldIt
]

\setsansfont{Roboto}[
  UprightFont = *-Regular,
  BoldFont = *-Bold,
  ItalicFont = *-Italic,
  BoldItalicFont = *-BoldItalic
]

% --- Typography ---
\usepackage[protrusion=true, expansion=true, final]{microtype}
\usepackage{setspace}
\setstretch{1.12}
\setlength{\parindent}{0pt}
\setlength{\parskip}{4pt}

% --- Colors ---
\usepackage[dvipsnames,svgnames,x11names,table]{xcolor}

\definecolor{PrimaryBlue}{HTML}{0066CC}
\definecolor{NoteBlue}{HTML}{3498DB}
\definecolor{WarningOrange}{HTML}{E67E22}
\definecolor{TipGreen}{HTML}{27AE60}
\definecolor{CodeBg}{HTML}{F5F5F5}
\definecolor{CodeFrame}{HTML}{CCCCCC}
\definecolor{DarkText}{HTML}{222222}
\definecolor{OfficialPurple}{HTML}{7B2D8E}
\definecolor{PractitionerTeal}{HTML}{008080}
\definecolor{WhyAmber}{HTML}{D4AC0D}
\definecolor{ExerciseIndigo}{HTML}{5B2C6F}
\definecolor{CostRed}{HTML}{CC3333}

% --- Code ---
\usepackage{minted}
\setminted{
  frame=lines,
  framesep=1mm,
  bgcolor=CodeBg,
  rulecolor=CodeFrame,
  fontsize=\footnotesize,
  breaklines=true,
  breakanywhere=false,
  linenos=false,
  xleftmargin=0pt,
  baselinestretch=1.05,
  autogobble=true
}
\usemintedstyle{friendly}

\usepackage[english]{babel}
\usepackage{csquotes}
\newcommand{\code}[1]{\texttt{#1}}

% --- Boxes ---
\usepackage[breakable,skins]{tcolorbox}

\tcbset{
  qs-box/.style={
    breakable,
    enhanced,
    arc=0pt,
    leftrule=3pt,
    rightrule=0pt,
    toprule=0pt,
    bottomrule=0pt,
    left=4pt,
    right=4pt,
    top=3pt,
    bottom=3pt,
    boxsep=0pt,
    toptitle=1pt,
    coltitle=DarkText,
    fonttitle=\bfseries\small
  }
}

\newtcolorbox{keybox}[1][Key Idea]{
  qs-box,
  colframe=PrimaryBlue,
  colback=PrimaryBlue!4,
  title=#1
}

\newtcolorbox{warnbox}[1][Common Mistake]{
  qs-box,
  colframe=WarningOrange,
  colback=WarningOrange!4,
  title=#1
}

\newtcolorbox{tipbox}[1][Tip]{
  qs-box,
  colframe=TipGreen,
  colback=TipGreen!4,
  title=#1
}

\newtcolorbox{whybox}[1][Why This Matters]{
  qs-box,
  colframe=WhyAmber,
  colback=WhyAmber!4,
  title=#1
}

% --- Tables ---
\usepackage{booktabs}
\usepackage{tabularx}
\usepackage{array}
\renewcommand{\arraystretch}{1.2}

% --- Lists ---
\usepackage[shortlabels]{enumitem}
\setlist{
  topsep=2pt,
  itemsep=1.5pt,
  parsep=0pt,
  partopsep=0pt,
  leftmargin=*,
  listparindent=0pt
}

% --- Graphics ---
\usepackage{tikz}
\usetikzlibrary{arrows.meta, positioning, shapes.geometric, calc}

% --- Hyperlinks ---
\usepackage[
  colorlinks=true,
  linkcolor=PrimaryBlue,
  urlcolor=PrimaryBlue,
  pdfusetitle,
  bookmarks=true,
  pdfstartview=FitH
]{hyperref}

% --- Headers ---
\usepackage[headsepline]{scrlayer-scrpage}
\pagestyle{scrheadings}
\ihead{\footnotesize Claude Code Quick Start}
\chead{}
\ohead{\footnotesize\textit{\headmark}}
\ifoot{} \cfoot{\pagemark} \ofoot{}

% --- Section styling ---
\addtokomafont{section}{\color{PrimaryBlue}\itshape}
\addtokomafont{subsection}{\color{PrimaryBlue!85!black}\itshape}

\RedeclareSectionCommand[
  beforeskip=0.8\baselineskip,
  afterskip=0.2\baselineskip
]{section}

\RedeclareSectionCommand[
  beforeskip=0.5\baselineskip,
  afterskip=0.15\baselineskip
]{subsection}

% =====================================================================
\begin{document}

\hypersetup{
  pdftitle={Claude Code Quick Start},
  pdfauthor={Brandon Behring},
  pdfsubject={Quick Start Guide for Claude Code},
  pdfkeywords={Claude Code, AI, LLM, TDD, quick start, 2026}
}

% --- Title ---
\begin{center}
  {\LARGE\bfseries\color{PrimaryBlue}Claude Code Quick Start}\\[4pt]
  {\large Your First 60 Minutes: From Setup to Verified Code}\\[6pt]
  {\normalsize Brandon Behring \quad|\quad February 2026}\\[2pt]
  {\footnotesize For data scientists and developers familiar with LLMs, new to Claude Code}
\end{center}

\vspace{0.3cm}
\hrule height 0.5pt
\vspace{0.4cm}

% =====================================================================
\section{The Mindset Shift (Minutes 0--5)}

You already know LLMs. Claude Code is not a chat interface---it is an
\textbf{agent loop} running in your terminal:

\begin{center}
\begin{tikzpicture}[
  node distance=1.4cm,
  box/.style={draw=PrimaryBlue, fill=PrimaryBlue!6, rounded corners=3pt,
    font=\small\bfseries, minimum width=1.6cm, minimum height=0.65cm, align=center},
  arr/.style={-{Stealth[length=5pt]}, thick, PrimaryBlue!70!black},
]
  \node[box] (prompt) {Prompt};
  \node[box, right=of prompt] (reason) {Reason};
  \node[box, right=of reason] (tool) {Tool};
  \node[box, right=of tool] (observe) {Observe};

  \draw[arr] (prompt) -- (reason);
  \draw[arr] (reason) -- (tool);
  \draw[arr] (tool) -- (observe);
  \draw[arr, rounded corners=8pt] (observe) -- ++(0,-0.8) -| (reason)
    node[pos=0.25, below, font=\tiny\itshape, text=DarkText] {repeat until done};
\end{tikzpicture}
\end{center}

Three properties define the system:

\begin{description}[leftmargin=2.5cm, style=nextline, font=\bfseries\color{PrimaryBlue!80!black}]
  \item[Context Window] A finite workspace. Every file read, tool result,
    and conversation turn consumes tokens. When it fills, early information fades.
    This is the single most important constraint to internalize.

  \item[Tool Use] Claude calls tools (\code{Read}, \code{Edit}, \code{Bash},
    \code{Grep}\ldots) instead of typing code. Each call costs context tokens
    and produces observable effects. Your job is to \emph{minimize unnecessary calls}
    through good configuration.

  \item[Configuration] Four layers shape behavior: \code{CLAUDE.md} (project knowledge),
    rules (conditional context), settings (permissions), and hooks (enforced automation).
    These are not suggestions---they are the operating system.
\end{description}

\begin{keybox}[The Core Insight]
Developers who treat Claude as a chatbot get chatbot results.
Developers who understand the agent loop---finite context, tool costs,
configuration layers---get a powerful engineering partner.
\end{keybox}

% =====================================================================
\section{Setup: CLAUDE.md and Deny Rules (Minutes 5--10)}

\code{CLAUDE.md} is a briefing document for a fast, capable, amnesiac teammate.
It is the \textbf{single most important file} in your configuration.

\begin{minted}{markdown}
# Demand Forecasting

## Build & Verify
- `pytest tests/` — run test suite
- `mypy src/` — type checking
- `ruff check src/` — lint

## Architecture
- src/features/ — feature engineering pipelines
- src/models/ — model training and evaluation
- src/data/ — data loading and validation
- notebooks/ — exploratory analysis
- tests/ — mirrors src/ structure

## Standards
- Type hints on all function signatures
- Run tests after every code change
- Never commit without passing lint + tests
\end{minted}

\begin{warnbox}
Do NOT write: \emph{``This is an ML project. We use scikit-learn. Tests are important.''}\\
DO write: \textbf{commands}, \textbf{locations}, and \textbf{rules}---not descriptions.
Claude needs to know how to \emph{work on} your project, not \emph{about} it.
\end{warnbox}

\textbf{Quick start}: Open your project, type \code{/init}.
Claude generates a starter CLAUDE.md from your project structure. Review, refine, commit.

\subsection{Protect Your Secrets (Day One)}

Create \code{.claude/settings.json}:

\begin{minted}{json}
{
  "permissions": {
    "deny": [
      "Read(.env*)", "Read(secrets/**)",
      "Read(**/*.pem)", "Read(**/*.key)"
    ]
  }
}
\end{minted}

Deny rules are absolute---cannot be overridden by prompts or prompt injection.

% =====================================================================
\section{Your First Verified Win (Minutes 10--25)}

\begin{whybox}
AI-generated code has a specific failure mode: it \textbf{looks correct}.
Clean syntax, reasonable names, plausible logic.
The bugs are subtle, not obvious. Verification is essential.
\end{whybox}

\subsection{The Pattern: Test First, Then Implement}

\begin{enumerate}
  \item \textbf{Describe the interface}: inputs, outputs, edge cases.
  \item \textbf{Claude writes tests} based on the interface.
  \item \textbf{Claude writes implementation} that passes the tests.
  \item \textbf{Tests run automatically} (via hooks or your prompt).
\end{enumerate}

\begin{minted}{text}
Prompt: "Create a FeatureValidator class:
- Takes a DataFrame and a schema dict
- Validates column types, value ranges, null counts
- Returns ValidationResult with errors list
- Raises ValueError if required columns are missing

Write tests first, then implementation."
\end{minted}

This works because Claude excels at test generation when given clear specs.
Tests then \emph{constrain} the implementation, preventing the
``looks correct but is subtly wrong'' failure mode.

\subsection{Verification Criteria in Every Prompt}

Always tell Claude how to check its own work:

\begin{minted}{text}
Add input validation to the build_features() function.
Verify: pytest tests/features/test_build.py passes,
and test with:
python -c 'from src.features import build_features;
  build_features(pd.DataFrame())'
(should raise ValueError, not silently return empty).
\end{minted}

\begin{keybox}[Anthropic's \#1 Practice]
``Give Claude a way to verify its own work'' is Anthropic's stated
single highest-leverage practice. It aligns exactly with what
practitioners discover independently.
\end{keybox}

\subsection{The 6-Layer Validation Architecture}

Layer defenses so no single failure mode goes undetected:

\begin{center}
\begin{tabular}{@{}clp{7.5cm}@{}}
\toprule
\textbf{\#} & \textbf{Layer} & \textbf{What It Catches} \\
\midrule
1 & Type Safety & Type errors at compile/lint time \\
2 & Input Validation & Precondition failures at function entry \\
3 & Unit Tests & Single-function correctness (80\%+ coverage) \\
4 & Integration Tests & Multi-component workflow failures \\
5 & End-to-End Tests & Complete user workflow failures \\
6 & Property-Based Tests & Edge cases you didn't think of \\
\bottomrule
\end{tabular}
\end{center}

Layers 1--2 are free to add today. Layers 3--4 are the baseline for production.

\subsection{Enforce with Hooks, Not Hope}

CLAUDE.md instructions are advisory---Claude may forget in long sessions.
Hooks are \textbf{enforced}:

\begin{minted}{json}
{
  "hooks": {
    "PostToolUse": [{
      "matcher": "Edit|Write",
      "hooks": [{
        "type": "command",
        "command": ".claude/hooks/lint-on-write.sh",
        "description": "Auto-lint written files"
      }]
    }],
    "PreToolUse": [{
      "matcher": "Bash",
      "hooks": [{
        "type": "command",
        "command": ".claude/hooks/pre-commit-gate.sh",
        "description": "Block commits with failing tests"
      }]
    }]
  }
}
\end{minted}

\textbf{Rule of thumb}: If a standard is non-negotiable, make it a hook.
If it is advisory, put it in CLAUDE.md.

% =====================================================================
\section{Plan Before You Edit (Minutes 25--35)}

For complex or unfamiliar tasks, separate exploration from implementation.
Press \code{Shift+Tab} twice to enter Plan Mode---Claude reads and analyzes
but cannot modify files.

\begin{enumerate}
  \item \textbf{Explore} (Plan Mode): ``Read \code{src/features/} and explain
    how temporal features are computed.''
  \item \textbf{Plan} (Plan Mode): ``I want to add recency features.
    What files change? Create a plan.''
  \item \textbf{Implement} (Normal Mode): ``Implement the plan.
    Write tests first, then code.''
  \item \textbf{Commit}: ``Commit with a descriptive message.''
\end{enumerate}

\begin{tipbox}[When to Skip]
Skip planning when the scope is clear and the fix is small.
If you can describe the diff in one sentence, just ask Claude to do it.
Plan when you are uncertain, the change spans multiple files,
or you are unfamiliar with the code.
\end{tipbox}

% =====================================================================
\section{The Session Loop (Minutes 35--45)}

Every productive session follows this cycle:

\begin{enumerate}
  \item \textbf{Prompt} --- describe what you want with verification criteria.
  \item \textbf{Observe} --- read the diff, check the tests, examine output.
  \item \textbf{Refine} --- give specific feedback (file, line, what to change).
  \item \textbf{Commit} --- one logical change per commit.
\end{enumerate}

\subsection{Context Hygiene (The Rules That Matter)}

\begin{description}[leftmargin=3.5cm, style=nextline, font=\ttfamily\bfseries]
  \item[/clear] Wipe context between unrelated tasks. Cost: 30 seconds.
  \item[/compact] Summarize, keeping key decisions. Use within long tasks.
  \item[Two-failure rule] After two failed corrections, \code{/clear} and write a
    better prompt. Three rounds of ``no, that's not what I meant'' means the
    context is dominated by failure patterns, not the task.
\end{description}

\begin{tipbox}
Watch the token counter. When it approaches 60--70\% of the window,
proactively \code{/compact} or \code{/clear}.
Context degradation is non-linear---there is a threshold where quality drops sharply.
\end{tipbox}

% =====================================================================
\section{Phase-Appropriate Rigor (Minutes 45--55)}

Not all code needs the same standards. Applying production rigor to a prototype
kills velocity. Applying prototype standards to production kills reliability.

\begin{center}
\begin{tabular}{@{}lp{3cm}p{3cm}p{4cm}@{}}
\toprule
\textbf{Phase} & \textbf{Testing} & \textbf{Code Quality} & \textbf{Transition Trigger} \\
\midrule
Exploration & Manual OK & Long functions OK & Hypothesis validated \\
Development & Unit + integration & Style enforced & Coverage >60\% \\
Production & Full 6-layer & Strict lint & Coverage >80\%, zero warnings \\
\bottomrule
\end{tabular}
\end{center}

Make phase transitions explicit in CLAUDE.md:

\begin{minted}{markdown}
## Phase: DEVELOPMENT (since Feb 15)
- Coverage target: 60%+
- Graduation criteria:
  - [ ] All feature pipelines have unit tests
  - [ ] No data leakage in CV pipeline
  When all checked: switch to PRODUCTION phase
\end{minted}

% =====================================================================
\section{The Six Traps (Reference)}

\begin{enumerate}[leftmargin=*, itemsep=4pt]
  \item \textbf{Context Overload} --- CLAUDE.md grows to 600+ lines,
    rules get ignored. \emph{Fix}: Keep under 300 lines.
    Move module-specific rules to \code{.claude/rules/} with path filters.

  \item \textbf{Kitchen Sink Session} --- Multiple unrelated tasks in one
    session. Quality degrades. \emph{Fix}: \code{/clear} between tasks.
    One logical task per session.

  \item \textbf{Over-Correcting} --- Four rounds of corrections, each adding noise.
    \emph{Fix}: Two-failure rule. \code{/clear}, write a better prompt.

  \item \textbf{Verification Gap} --- AI code accepted without tests.
    Subtle bugs surface weeks later. \emph{Fix}: Test-first. Hooks enforce it.

  \item \textbf{Permanent Prototype} --- ``We'll add tests later'' becomes never.
    \emph{Fix}: Explicit phase transitions in CLAUDE.md with dates and criteria.

  \item \textbf{Infinite Exploration} --- ``Investigate how auth works'' without scoping.
    Claude reads 40 files, filling context with exploration results.
    \emph{Fix}: Scope investigations narrowly. Use subagents for research
    so exploration runs in isolated context.
\end{enumerate}

\begin{keybox}[The Common Thread]
Every trap comes from treating Claude as infinitely capable.
Context is finite. Attention degrades. Verification is essential.
Working \emph{within} these constraints---not ignoring them---is
what separates productive from frustrating usage.
When stuck: \code{Esc} to stop, \code{Esc + Esc} to rewind,
\code{/clear} to start fresh.
\end{keybox}

% =====================================================================
\section{Quick Reference Card}

\subsection{Essential Commands}

\begin{center}
\begin{tabular}{@{}lp{8cm}@{}}
\toprule
\textbf{Command / Shortcut} & \textbf{What It Does} \\
\midrule
\code{/init} & Generate starter CLAUDE.md from project \\
\code{/clear} & Clear context, start fresh \\
\code{/compact} & Summarize, preserve key decisions \\
\code{/context} & Visualize current context usage \\
\code{/rewind} & Restore conversation, code, or both to a checkpoint \\
\code{Esc} & Stop Claude mid-action (context preserved) \\
\code{Esc + Esc} & Open rewind menu (restore or summarize from checkpoint) \\
\code{Shift+Tab} & Cycle modes (Normal $\rightarrow$ Auto-Accept $\rightarrow$ Plan) \\
\code{Alt+T} & Toggle extended thinking \\
\code{--continue} & Resume most recent session \\
\code{--resume} & Pick from previous sessions \\
\bottomrule
\end{tabular}
\end{center}

\subsection{The One-Page Workflow}

\begin{center}
\begin{tikzpicture}[
  node distance=0.4cm and 1.0cm,
  wfstep/.style={draw=PrimaryBlue, fill=PrimaryBlue!6, rounded corners=3pt,
    font=\small, minimum width=5.5cm, minimum height=0.8cm, align=left},
  num/.style={font=\bfseries\small\color{PrimaryBlue}},
  arr/.style={-{Stealth[length=5pt]}, thick, PrimaryBlue!50},
]
  \node[wfstep] (s1) {\code{/init} \textrightarrow{} create CLAUDE.md + deny rules};
  \node[wfstep, below=of s1] (s2) {Add verification criteria + hook enforcement};
  \node[wfstep, below=of s2] (s3) {Plan Mode if complex; direct prompt if simple};
  \node[wfstep, below=of s3] (s4) {Prompt with specs + ``write tests first''};
  \node[wfstep, below=of s4] (s5) {Observe output, check diffs and tests};
  \node[wfstep, below=of s5] (s6) {Refine (\code{Esc} to stop, feedback) or \code{/clear} + retry};
  \node[wfstep, below=of s6] (s7) {Commit one logical change};

  \node[num, left=0.3cm of s1] {1};
  \node[num, left=0.3cm of s2] {2};
  \node[num, left=0.3cm of s3] {3};
  \node[num, left=0.3cm of s4] {4};
  \node[num, left=0.3cm of s5] {5};
  \node[num, left=0.3cm of s6] {6};
  \node[num, left=0.3cm of s7] {7};

  \draw[arr] (s1) -- (s2);
  \draw[arr] (s2) -- (s3);
  \draw[arr] (s3) -- (s4);
  \draw[arr] (s4) -- (s5);
  \draw[arr] (s5) -- (s6);
  \draw[arr] (s6) -- (s7);

  % Loop back
  \draw[arr, rounded corners=6pt] (s7.east) -- ++(0.6,0) |- (s3.east)
    node[pos=0.25, right, font=\tiny\itshape, text=DarkText] {next task};
\end{tikzpicture}
\end{center}

\vfill

\begin{center}
\footnotesize
\textit{Extracted from ``Best Practices for Using Claude'' (v2.1+).}\\
\textit{Full handbook: 16 chapters covering foundations through enterprise deployment.}\\[4pt]
Licensed under CC~BY~4.0 \quad|\quad Content verified against Anthropic docs, February 2026
\end{center}

\end{document}
